\documentclass[conference]{IEEEtran}
\IEEEoverridecommandlockouts

\usepackage{cite}
\usepackage{amsmath,amssymb,amsfonts}
\usepackage{graphicx}
\usepackage{textcomp}
\usepackage{xcolor}
\usepackage{hyperref}
\usepackage{booktabs}
\usepackage{array}
\usepackage{url}
\usepackage[utf8]{inputenc}
\usepackage{listings}
\usepackage{balance}
\usepackage{dblfloatfix}
\usepackage{float}
\usepackage{caption}
\usepackage[section]{placeins}  % \FloatBarrier keeps floats within their section

\lstset{
  basicstyle=\ttfamily\scriptsize,
  breaklines=true,
  frame=single,
  columns=fullflexible,
}

\def\BibTeX{{\rm B\kern-.05em{\sc i\kern-.025em b}\kern-.08em
    T\kern-.1667em\lower.7ex\hbox{E}\kern-.125emX}}

\begin{document}

\title{GitHub Technology Trend Analysis:\\
A Large-Scale Big Data Pipeline for Developer Ecosystem Intelligence\\
{\footnotesize Phase 2 -- Midterm Report}}

\author{
  \IEEEauthorblockN{Selin Bardakc\i}
  \IEEEauthorblockA{
    \textit{Dept.\ of Computer Engineering}\\
    \textit{Gebze Technical University}\\
    sbardakci@gtu.edu.tr
  }
  \and
  \IEEEauthorblockN{Bengisu Duru G\"{o}ksu}
  \IEEEauthorblockA{
    \textit{Dept.\ of Computer Engineering}\\
    \textit{Gebze Technical University}\\
    bgoksu2022@gtu.edu.tr
  }
  \and
  \IEEEauthorblockN{Zeynep Yi\u{g}it}
  \IEEEauthorblockA{
    \textit{Dept.\ of Computer Engineering}\\
    \textit{Gebze Technical University}\\
    zyigit@gtu.edu.tr
  }
}

\maketitle

% ─────────────────────────────────────────────────────────────────────────────
\begin{abstract}
We present a scalable big data pipeline that ingests, processes, and analyses
over 22.9 million GitHub repositories to surface actionable insights about
programming language adoption, technology migration patterns, and developer
collaboration networks. The pipeline integrates two public data sources — the
HuggingFace \textit{github-repos-metadata-40M} dataset and the GitHub Archive
event stream — processes them through six sequential Apache Spark jobs on
Google Cloud Dataproc, and produces five analytical products: monthly language
trend metrics (542 languages, 126 months), 24-month Holt-Winters ETS forecasts,
a technology migration graph, developer influence PageRank scores, and Leiden
community detection over 25,000 developers. This report documents all
engineering decisions, design trade-offs, failed approaches and their
replacements, and preliminary quantitative results across all five analyses.
\end{abstract}

\begin{IEEEkeywords}
big data, GitHub, technology trends, Apache Spark, time-series forecasting,
Holt-Winters ETS, graph analytics, PageRank, Leiden community detection,
KMeans clustering, GCP Dataproc
\end{IEEEkeywords}

\begin{table}[htbp]
\centering\small
\begin{tabular}{p{0.9\columnwidth}}
\hline
\textbf{Code \& Data Availability}\\
\hline
\vspace{1pt}
Full source code, Spark jobs, and notebooks are publicly available:\\[3pt]
\url{https://github.com/selin-bardakci/BigData-Github_Analysis}
\vspace{3pt}\\
\hline
\end{tabular}
\end{table}

% ─────────────────────────────────────────────────────────────────────────────
\section{Revised Problem Statement}

\subsection{Original Proposal}

The Phase 1 proposal aimed to characterise programming language popularity
trends using the static \texttt{bigquery-public-data.github\_repos} snapshot
(data frozen circa 2017) and apply ARIMA(2,1,2) forecasting with a basic
co-authorship graph computed inside BigQuery.

\subsection{Fundamental Data Source Problem}

Upon deeper investigation, the original BigQuery dataset presented two
critical limitations. First, the language snapshot is \emph{point-in-time}:
it records the byte count of each language in each repository as of a single
crawl date, not as a time series. This means any time-series analysis would
have been driven entirely by which repositories happened to be active during
the crawl window, producing a cross-sectional artefact rather than genuine
adoption dynamics. Second, the dataset covers only repositories that existed
by 2017, systematically excluding the 2018--2025 rise of TypeScript, Rust,
Go, and Kotlin — the very languages whose growth trajectories are most
informative for forecasting.

\subsection{Revised Data Architecture}

Both sources were replaced:

\textbf{Q1 (Language labels):} Replaced the 2017 snapshot with the HuggingFace
\texttt{ibragim-bad/github-repos-metadata-40M} dataset — 22.9 million
repositories labelled with GitHub's Linguist tool, spanning creation dates
through 23 July 2025, with an explicit \texttt{created\_at} timestamp per
repository. This dataset provides the same language labels as the original
but extends coverage by eight years and includes temporal metadata that
enables the pre-birth filter described in Section~\ref{sec:preprocess}.

\textbf{Q2 (Commit activity):} Replaced the frozen commits table with
\texttt{githubarchive.month.*} PushEvent counts — a live BigQuery dataset
covering monthly commit activity from 2015 to present. This provides a
true time series of repository activity with which Q1 language labels can
be cross-joined.

\textbf{Q3 (Migration signals):} Extracted technology migration intent from
GH Archive commit messages using SQL pattern matching on verb-noun pairs
(e.g., ``migrating from Flask to FastAPI'').

\textbf{Q4 (Developer activity):} Extracted per-developer, per-repository,
per-year commit aggregates from GH Archive PushEvents, enabling construction
of a developer co-contribution graph.

\subsection{Revised Research Questions}

\begin{enumerate}
  \item Which programming languages show statistically significant growth or
        decline across 2015--2025, and how do they cluster by adoption trajectory?
  \item What are the 24-month Holt-Winters forecasts for the top-30 languages,
        and which are projected to gain or lose ecosystem share?
  \item Which technologies act as migration destinations (high in-flow) vs.\
        migration sources (high out-flow) in the developer ecosystem?
  \item How is influence distributed among active open-source developers,
        and do natural collaboration communities align with language ecosystems?
\end{enumerate}

% ─────────────────────────────────────────────────────────────────────────────
\section{Related Work}

Dyer et al.~\cite{dyer2013boa} built the Boa infrastructure for large-scale
GitHub mining, demonstrating that distributed storage and compute are
necessary for repository-scale analysis. Our pipeline adopts the same
architectural pattern (GCS + Spark) but focuses on temporal adoption dynamics
rather than static code properties.

Borges et al.~\cite{borges2016popularity} studied what drives repository
popularity (stars, forks, commits), finding that language choice is a
significant predictor. Our work quantifies language-level trend behaviour
at ecosystem scale rather than per-repository.

For forecasting, Gardner~\cite{gardner1985exponential} established the
theoretical basis for exponential smoothing with damped trend, and the M4
competition~\cite{makridakis2018m4} empirically demonstrated that Holt-Winters
ETS outperforms ARIMA on seasonal data in the short- to medium-term horizon.
We selected Holt-Winters over ARIMA for exactly this reason (Section~\ref{sec:d2}).
Hyndman and Athanasopoulos~\cite{hyndman2018forecasting} provide the
state-space ETS formulation implemented in \texttt{statsmodels}.

Community detection follows the Leiden algorithm~\cite{traag2019leiden},
which corrects a well-known defect in the Louvain method: Louvain can produce
internally disconnected communities because it only optimises modularity
locally. Leiden adds a refinement step that guarantees all communities are
well-connected before the final partition is returned. We specifically chose
Leiden over Louvain for the developer graph because isolated sub-clusters
(small teams within a larger project) would be incorrectly merged by Louvain.
Developer influence is measured with PageRank~\cite{brin1998pagerank}
($\alpha = 0.85$) rather than degree centrality, as degree is easily gamed
by contributing trivially to many repositories.

% ─────────────────────────────────────────────────────────────────────────────
\section{System Architecture \& Design}

\subsection{Overall Pipeline}

The pipeline consists of four sequential stages:

\begin{enumerate}
  \item \textbf{Ingestion:} Data is pulled from external sources into GCS
        raw storage as Parquet + Snappy shards.
  \item \textbf{Preprocessing (Spark):} Six Spark jobs clean, join, aggregate,
        cluster, forecast, and graph-analyse the raw data, writing processed
        Parquet to GCS.
  \item \textbf{Analysis (Jupyter):} Five domain notebooks and one Final
        Report notebook read processed GCS outputs and produce all
        visualisations.
  \item \textbf{Reporting:} Figures exported from notebooks are compiled into
        this report.
\end{enumerate}

Fig.~\ref{fig:pipeline} shows the complete data flow.

%% ── PIPELINE FIGURE ──────────────────────────────────────────────────────────
%% Draw this in draw.io / Lucidchart / TikZ using the following 3-tier layout:
%%
%% TIER 1 — SOURCES (left column, 4 boxes):
%%   [HuggingFace 40M]  →  [GCS raw/languages/  459 shards, 22.9M rows]
%%   [GH Archive BQ: PushEvents] → [GCS raw/commits/]
%%   [GH Archive BQ: Migrations] → [GCS raw/migrations/]
%%   [GH Archive BQ: Developers] → [GCS raw/developers/]
%%
%% TIER 2 — DATAPROC / SPARK (centre, large rounded rectangle):
%%   Sequential chain:  d1_preprocess → d1_aggregate → d1_features_cluster
%%                                                   ↘ d2_forecasting
%%   Parallel branches: d3_migration_graph
%%                      d4d5_developer_graph
%%
%% TIER 3 — OUTPUTS (right column):
%%   D1: processed/d1_monthly/  processed/d1_clusters/
%%   D2: processed/d2_forecasts/
%%   D3: processed/d3_edges/  processed/d3_nodes/
%%   D4: processed/d4_pagerank/
%%   D5: processed/d5_communities/
%%   ↓ (all) → [Jupyter Notebooks: D1–D5 + Final_Report]
%%
%% Style: blue=GCS boxes, orange=Spark jobs, green=notebooks.
%% Export as PDF → figures/pipeline.pdf, then use \includegraphics.
%% ─────────────────────────────────────────────────────────────────────────────

\begin{figure*}[!t]
  \centering
  \includegraphics[width=\linewidth]{pipeline}
  \caption{End-to-end data pipeline. Data flows from HuggingFace and GitHub
    Archive through GCS raw storage, six Spark jobs on Dataproc, and into
    five analytical GCS outputs consumed by Jupyter notebooks.}
  \label{fig:pipeline}
\end{figure*}

\subsection{Storage: Parquet + Snappy on GCS}

All inter-stage data is stored as Apache Parquet with Snappy compression.
This choice was made for three reasons:

\textit{(1) Columnar reads.} Spark jobs typically access 2--4 columns from
wide schemas. Parquet's columnar layout means only the needed columns are
read from disk, reducing I/O by $\sim$5$\times$ compared to CSV.

\textit{(2) Predicate pushdown.} Parquet row-group statistics allow Spark to
skip row groups entirely when a filter predicate (e.g., \texttt{year\_month
<= "2025-06"}) cannot match. This is used extensively in \texttt{d1\_preprocess}.

\textit{(3) Splittability.} Each Parquet file is independently readable by
separate Spark executors, enabling fine-grained parallelism. The
HuggingFace dataset was deliberately sharded into 459 files during ingestion
to allow 459-way parallelism in the first Spark read.

Snappy was chosen over GZIP because Snappy decompresses 5$\times$ faster at
the cost of $\sim$20\% larger files. For compute-bound Spark jobs, decompression
speed dominates total job time.

\subsection{Compute: Google Cloud Dataproc}

Dataproc was chosen over raw GCE instances because it provides managed Spark
cluster lifecycle (auto-start, auto-teardown), built-in YARN resource manager,
and native GCS connector without needing HDFS.

\textbf{Zone selection issue.} The initial cluster was configured for
\texttt{us-central1}. Two consecutive provisioning attempts failed with
``Out of capacity'' errors on zones \texttt{us-central1-a} and
\texttt{us-central1-f}. The cluster was moved to \texttt{us-east1} with
auto-zone selection (\texttt{-{}-zone=""}) to allow GCP to pick a zone
with available capacity.

\textbf{Job orchestration.} An early implementation submitted all six Spark
jobs simultaneously with the \texttt{-{}-async} flag. This caused a race
condition: \texttt{d1\_aggregate} began reading \texttt{d1\_preprocessed/}
before \texttt{d1\_preprocess} had finished writing it, resulting in corrupt
Parquet footer errors. The orchestration script
(\texttt{scripts/run\_all\_spark\_jobs.sh}) was rewritten to submit jobs
sequentially, polling for completion via \texttt{gcloud dataproc jobs wait}
before submitting the next dependent job.

\textbf{Driver-side library installation.} Dataproc base images do not
include \texttt{statsmodels}, \texttt{python-igraph}, or \texttt{leidenalg}.
Each Spark job that requires a non-standard library uses a helper
\texttt{\_ensure\_module()} function that attempts \texttt{importlib.import\_module}
first and falls back to \texttt{pip install} at runtime if the module is
absent. This avoids the need for a custom Dataproc image while keeping the
job scripts self-contained.

\subsection{Data Quality: The 2025-06 Cutoff Decision}

The HuggingFace dataset covers repositories created through 23 July 2025.
When the D1 inner join (\texttt{languages} $\bowtie$ \texttt{commits}) is
evaluated for months after July 2025, repositories created after the crawl
date have no language label and fall out of the join. This produces an
artefact: the monthly repository count drops by 49\% between June and August
2025 for languages like JavaScript (153k~$\rightarrow$~79k in one month),
mimicking a catastrophic ecosystem collapse that did not occur.

All D1 and D2 data are therefore hard-capped at \texttt{year\_month <= "2025-06"},
giving 10.5 years (126 months) of clean, complete monthly observations.
The D4/D5 developer graph applies a complementary filter (\texttt{year <= "2025"})
to exclude partial-year 2026 data from GH Archive.

% ─────────────────────────────────────────────────────────────────────────────
\section{Implementation: Spark Jobs}
\label{sec:implementation}

\subsection{D1 Preprocessing (\texttt{d1\_preprocess.py})}
\label{sec:preprocess}

\textbf{Input:} \texttt{raw/languages/} (Q1: repo\_name, language, bytes,
created\_at) and \texttt{raw/commits/} (Q2: repo\_name, year\_month,
commit\_count).

\textbf{Language filtering.} A static exclusion set removes 19 categories
that are not programming languages: Markdown, JSON, YAML, XML, HTML, CSS,
Dockerfile, CMake, Shell, Makefile, and similar config/markup formats.
These appear frequently in repositories but carry no signal about programming
language adoption.

\textbf{Inner join.} A Spark inner join on \texttt{repo\_name} produces one
row per (repo, language, year\_month). Because Q1 is a flat table (one row
per language per repo, not per month), joining Q2's monthly rows cross-assigns
each language's static snapshot onto every month the repo had commits.

\textbf{Pre-birth filter.} A critical quality filter drops any
(repo, year\_month) row where \texttt{year\_month < created\_month}.
Without this filter, a repository created in 2022 would have its
TypeScript label applied retroactively to 2015--2021 months based on Q2
commit history from that period (possibly inherited commits or forks).
This was identified as a key cause of an artificial spike in TypeScript
counts before 2019. Rows where \texttt{created\_at} is NULL are conservatively
retained (no filter applied).

\textbf{Output partitioning.} The preprocessed data is written with
\texttt{.partitionBy("language")} and 200 partitions to allow downstream
Spark jobs to read only the language partition they need without a full
dataset scan.

\subsection{D1 Aggregation (\texttt{d1\_aggregate.py})}

Groups the preprocessed data by (language, year\_month) to produce:
\texttt{repo\_count} (distinct repositories), \texttt{commit\_count} (total
commits), and \texttt{total\_bytes} (sum of bytes, set to 0 because the HF
dataset does not provide per-language byte counts). Languages with fewer than
12 months of coverage are filtered, ensuring all downstream models have at
least one full seasonal cycle of data.

\textbf{Result:} 53,630 rows covering 542 languages over 126 months.

\subsection{D1 Feature Engineering \& Clustering (\texttt{d1\_features\_cluster.py})}

\textbf{Feature design.} Five features were designed to capture distinct
aspects of a language's adoption trajectory:

\begin{enumerate}
  \item \texttt{avg\_repo\_growth}: Mean month-over-month percentage change
        in \texttt{repo\_count}. This is the primary adoption growth signal.
  \item \texttt{commit\_volatility}: Coefficient of variation of
        \texttt{commit\_count} ($\sigma / \mu$). High volatility indicates
        event-driven spikes (e.g., conference releases, library launches).
  \item \texttt{repo\_range\_ratio}: $\max(\text{repo\_count}) /
        \min(\text{repo\_count})$ over the full history. A ratio near 1
        indicates a stable, mature language; high ratios indicate languages
        still growing (or declining) rapidly.
  \item \texttt{avg\_repo\_count}: Mean monthly repository count. Encodes
        the absolute scale of adoption (distinguishing Python from Nim).
  \item \texttt{avg\_commit\_count}: Mean monthly commit count. Separates
        passive code repositories from actively maintained projects.
\end{enumerate}

\textbf{Bytes removed.} An initial feature set included bytes-based
features (\texttt{avg\_bytes\_growth}, \texttt{bytes\_range\_ratio}).
These were discarded because Q1 is a static snapshot: the bytes column
records the size of each repository at the crawl date and is cross-joined
unchanged onto every month. Month-over-month bytes growth is therefore zero
by construction, contributing pure noise to the clustering.

\textbf{Scaling.} The 5-D feature vectors are normalised with
\texttt{StandardScaler(withMean=True, withStd=True)} before KMeans to
prevent the high-magnitude \texttt{avg\_repo\_count} dimension from
dominating the Euclidean distance calculation.

\textbf{Clustering.} KMeans is run via Spark MLlib ($k=6$, \texttt{seed=42},
\texttt{maxIter=50}). $k=6$ was selected because it produces the highest
Silhouette score across $k \in \{3, 4, 5, 6, 7, 8\}$ and aligns with
natural language trajectory archetypes: fast-growing modern languages
(TypeScript, Rust, Go), dominant incumbents (JavaScript, Python, Java),
stable systems languages (C, C++), declining legacy languages (Ruby,
Objective-C), niche/domain-specific (R, Haskell, Julia), and long-tail
obscure languages ($\sim$430 languages with minimal repository presence).

\subsection{D2 Forecasting (\texttt{d2\_forecasting.py})}
\label{sec:d2}

\subsubsection{Initial Approach: ARIMA(2,1,2)}

The original forecasting model was ARIMA(2,1,2) from
\texttt{statsmodels.tsa.arima.model}. ARIMA differencing ($d=1$) removes
a linear trend, and the AR(2) and MA(2) terms model short-term autocorrelation
in the differenced series. This worked as a structural baseline but
produced qualitatively wrong forecasts: TypeScript had grown 250\% over
five years, yet ARIMA projected near-flat continuation at the most
recently observed level ($\sim$166k repos). Rust had grown 1,250\% over
10 years; ARIMA forecast near-zero growth.

\textbf{Root cause.} ARIMA extrapolates only short-term momentum. After
first differencing, the AR terms capture ``how much the growth rate
changed last month'' — not the 10-year secular trend. For a language that
has been growing 2\% per month consistently, ARIMA sees a stable differenced
series and forecasts continued 2\% growth only for a few steps before mean-reverting
to zero growth. The long-horizon forecast collapses to the mean.

\subsubsection{Replacement: Holt-Winters Exponential Smoothing}

ARIMA was replaced with Holt-Winters ETS (Error-Trend-Seasonal) from
\texttt{statsmodels.tsa.\allowbreak holtwinters.\allowbreak ExponentialSmoothing} using:

\begin{itemize}
  \item \texttt{trend="add"} — explicit additive trend state that
        accumulates directional momentum over the full history.
  \item \texttt{damped\_trend=True} — dampens the trend toward zero as
        the horizon extends, preventing the unrealistic straight-line
        extrapolation that undamped additive trend produces for mature
        languages (which slow as they reach market saturation).
  \item \texttt{seasonal="add"} with \texttt{seasonal\_periods=12} —
        captures the yearly GitHub activity cycle (January slowdown,
        summer peak) that is visible in all major languages.
  \item \texttt{initialization\_method="estimated"} — fits initial
        state values jointly with the smoothing parameters via MLE,
        rather than using a fixed heuristic.
\end{itemize}

A minimum history requirement of 36 months (3 full seasonal cycles) was
enforced so that Holt-Winters can reliably estimate the 12-month seasonal
component. All 30 top languages met this requirement (each has 126 months
of history). Gap-filling with linear interpolation handles the rare cases
where a language has no repositories in an isolated month; values are
clipped to a minimum of 1.0 before fitting to prevent issues with
log-transform variants of ETS.

\subsubsection{Prophet: Attempted and Abandoned}

Prophet (Meta/Facebook) was investigated as an alternative to both ARIMA
and Holt-Winters because it supports holiday effects, automatic seasonality
detection, and provides uncertainty intervals. Two blockers prevented its
use:

\textit{(1) Dataproc compatibility.} Prophet requires CmdStan (a C++ Stan
backend) to be compiled at runtime. On Dataproc base images, the Stan C++
compilation consistently timed out or produced build errors due to restricted
compiler flags. Attempting to pre-install CmdStan via \texttt{pip install
cmdstanpy} followed by \texttt{cmdstanpy.install\_cmdstan()} hung for
over six hours without completing.

\textit{(2) Import-time backend initialisation.} Even when installed in the
local virtual environment (not on Dataproc), Prophet 1.1.5 raises a
\texttt{stan\_backend} error if the CmdStan binary path is not set before
the first \texttt{import prophet}. The CmdStan setup cell must precede all
Prophet imports, but Jupyter's kernel resets this ordering when cells are
re-run out of sequence.

The schema retains \texttt{prophet\_yhat}, \texttt{prophet\_lower}, and
\texttt{prophet\_upper} columns (filled with NaN) for forward compatibility,
in case Prophet is enabled in a future controlled environment.

\subsection{D3 Migration Graph (\texttt{d3\_migration\_graph.py})}

\textbf{Signal extraction.} Q3 migration signals were extracted from GH
Archive commit messages using SQL pattern matching on phrases of the form
``migrated from X to Y'', ``switching from X to Y'', ``moving from X to Y'',
etc. The raw signal table contains (from\_tech, to\_tech, year,
migration\_count) tuples.

\textbf{Technology allowlist.} Raw extraction produces many false positives:
natural language words (``from'', ``to'', ``go'') are misinterpreted as
technology names, and abbreviations are ambiguous. A 70-technology allowlist
covering programming languages, frameworks, databases, CI/CD tools, and
infrastructure technologies was constructed. Both endpoints of each edge
must appear in the allowlist; any edge touching an unlisted technology is
discarded.

\textbf{Minimum edge weight.} Edges with fewer than 3 total migration signals
across all years are dropped (\texttt{MIN\_EDGE\_WEIGHT = 3}). This eliminates
single-mention noise while preserving weak-but-real migration patterns.

\textbf{Graph construction.} The filtered edges are collected to the Spark
driver (the migration graph is small after filtering: 16 nodes, 11 edges)
and loaded into a NetworkX directed graph. Per-node PageRank, betweenness
centrality, in-weight, out-weight, and net flow are computed on the driver
and written as Parquet.

\textbf{Limitation.} The Q3 dataset is sparse: only 11 edges survive the
quality filters. GH Archive commit messages are largely free-form and
rarely contain structured migration intent. Increasing coverage requires
either a lower allowlist threshold (risking more false positives) or an
NLP-based classifier applied to raw commit text — a direction reserved for
Phase 3.

\subsection{D4/D5 Developer Graph (\texttt{d4d5\_developer\_graph.py})}

\subsubsection{Privacy: SHA-256 Email Hashing}

Developer emails from GH Archive are irreversibly hashed with SHA-256
immediately upon DataFrame creation — before any Spark transformation,
any GCS write, or any log output. The raw email column is dropped from
the DataFrame immediately after the UDF assigns the hashed \texttt{dev\_id}.
No raw email appears in any processed output or log entry.

\subsubsection{Developer Selection: Distinct Repos vs.\ Total Commits}

The initial implementation selected the top developers by total commit
count. This produced a bot-dominated node set: CI automation accounts
that push millions of commits to a single monorepo were ranked at the top.
The resulting graph had 83\% isolated nodes and a mean degree of 0.6,
making PageRank and community detection meaningless.

The selection criterion was changed to \emph{distinct repository count}
(top 25,000 developers by number of unique repositories they contributed
to), with two additional filters:

\begin{itemize}
  \item \texttt{MAX\_REPOS\_PER\_DEV = 2000}: Developers contributing to
        more than 2,000 distinct repositories are almost certainly automated
        fork-and-commit bots (e.g., package mirror scripts).
  \item \texttt{MIN\_DISTINCT\_REPOS = 2}: Requires at least two distinct
        repositories to filter one-shot, drive-by commits.
\end{itemize}

After applying these filters, mean degree rose from 0.6 to 7.7, the
isolated node fraction dropped from 83\% to 38.5\%, and the graph became
suitable for graph analytics.

\subsubsection{Megarepo Exclusion}

The co-contribution self-join (``two developers share an edge if they both
committed to the same repository in the same year'') has an $O(N^2)$
complexity per repository, where $N$ is the number of contributors. For
megarepos like the Linux kernel or TensorFlow, $N > 10,000$, making the
pairwise join produce 50+ million edge rows from a single repository.

A pre-join filter counts the number of distinct developers per repository
and excludes repositories with more than \texttt{MAX\_DEVS\_PER\_REPO = 200}
contributors from the self-join. These megarepos contribute edges between
\emph{all} their contributors, which would dominate and homogenise the
graph. Excluding them preserves the graph's ability to capture smaller,
specialised collaboration clusters.

\subsubsection{HITS vs.\ Betweenness Centrality}

The initial D4 design computed HITS (Hyperlink-Induced Topic Search) hub
and authority scores alongside PageRank. HITS was intended to identify
``hub'' developers who link to many authorities and ``authority'' developers
who are linked from many hubs. After implementation, it was observed that
hub\_score $\equiv$ authority\_score for all nodes — a mathematical
inevitability on undirected graphs, where every edge is bidirectional and
there is no distinction between ``pointing to'' and ``being pointed at''.

HITS was replaced with betweenness centrality, which meaningfully measures
bridge nodes in an undirected co-contributor graph: a developer with high
betweenness lies on many shortest paths between other developer pairs,
indicating they connect otherwise-disjoint collaboration clusters.

\textbf{igraph API issue.} The \texttt{igraph.Graph.betweenness()} method
does not accept a \texttt{normalized=True} keyword argument (despite some
documentation suggesting otherwise). The call raised a \texttt{TypeError:
unexpected keyword argument}. Normalisation was applied manually:
$b_{\text{norm}}(v) = b(v) \;/\; \tfrac{(n-1)(n-2)}{2}$, where $n$ is the
number of vertices.

\subsubsection{Community Detection: Leiden over Louvain}

Leiden~\cite{traag2019leiden} was selected over Louvain for two reasons:

\textit{(1) Well-connectedness guarantee.} Louvain can produce internally
disconnected communities by merging sub-clusters that share a strong external
link. Leiden adds a refinement phase after each modularity-optimisation step
that splits such internally-disconnected communities before they can persist.
For a developer graph where sub-community identity is important (teams
within a project), this guarantee is critical.

\textit{(2) Reproducibility.} Leiden with a fixed seed (\texttt{seed=42})
produces deterministic partitions across runs. Louvain's community assignment
is non-deterministic in its original form.

% ─────────────────────────────────────────────────────────────────────────────
% ─────────────────────────────────────────────────────────────────────────────
\section{Preliminary Results}

\subsection{D1: Language Trend Analysis}

The monthly aggregation covers \textbf{542 languages} over
\textbf{126 months} (January 2015 -- June 2025), with 53,630 total rows.

The top-15 languages by mean monthly repository count are: JavaScript,
Python, TypeScript, Java, PHP, C++, C\#, Go, Ruby, C, Jupyter Notebook,
Rust, Vue, Kotlin, and Swift. Table~\ref{tab:top10} shows representative
statistics for the top-10.

\begin{table}[h]
\caption{Top-10 Languages: Mean Monthly Repository Count (2015--2025)}
\label{tab:top10}
\centering
\begin{tabular}{clr}
\toprule
\textbf{Rank} & \textbf{Language} & \textbf{Mean repos/month} \\
\midrule
1  & JavaScript      & $\sim$133,000 \\
2  & Python          & $\sim$96,000  \\
3  & TypeScript      & $\sim$47,000  \\
4  & Java            & $\sim$72,000  \\
5  & PHP             & $\sim$38,000  \\
6  & C++             & $\sim$31,000  \\
7  & C\#             & $\sim$28,000  \\
8  & Go              & $\sim$21,000  \\
9  & Ruby            & $\sim$24,000  \\
10 & C               & $\sim$20,000  \\
\bottomrule
\end{tabular}
\end{table}

\textbf{Note:} Values are approximate; exact figures are visible in the D1
notebook output. TypeScript's mean is lower than Java's despite faster
growth because TypeScript's growth is concentrated in 2020--2025 while
Java's mean is elevated by its large 2015--2019 base.

KMeans ($k=6$) partitions the 542 languages into trajectory archetypes.
Fig.~\ref{fig:d1_cluster} shows the clustering results; representative
members of each cluster are:

\begin{itemize}
  \item \textbf{Cluster: Fast-growing modern} — TypeScript, Rust, Go, Kotlin,
        Dart, HCL. High \texttt{avg\_repo\_growth}, high
        \texttt{repo\_range\_ratio}.
  \item \textbf{Cluster: Dominant incumbents} — JavaScript, Python, Java.
        High \texttt{avg\_repo\_count}, moderate growth rate.
  \item \textbf{Cluster: Stable systems languages} — C, C++, Objective-C.
        High \texttt{avg\_commit\_count}, low growth.
  \item \textbf{Cluster: Declining legacy} — Ruby, PHP, Perl, CoffeeScript.
        Negative \texttt{avg\_repo\_growth}, decreasing
        \texttt{repo\_range\_ratio}.
  \item \textbf{Cluster: Niche/academic} — R, Haskell, Julia, Erlang, OCaml.
        Low \texttt{avg\_repo\_count}, high \texttt{commit\_volatility}.
  \item \textbf{Cluster: Long-tail} — $\sim$430 languages. Near-zero
        \texttt{avg\_repo\_count} and \texttt{avg\_commit\_count}.
\end{itemize}

\begin{figure}[htbp]
  \centering
  \includegraphics[width=\columnwidth]{d1_cluster_pca}
  \caption{KMeans clustering of 542 languages projected onto the first two
    PCA components. Each point is one language; colour encodes cluster
    membership. JavaScript and Python (Cluster 1, orange) are extreme
    outliers on PC1, reflecting their uniquely large repository counts.}
  \label{fig:d1_cluster}
\end{figure}

\begin{figure}[htbp]
  \centering
  \includegraphics[width=\columnwidth]{d1_timeseries}
  \caption{Monthly repository count for the top-10 languages (January 2015 --
    June 2025). JavaScript dominates in absolute scale; TypeScript shows the
    steepest relative growth, rising from near-zero in 2015 to the third
    largest language by 2025.}
  \label{fig:d1}
\end{figure}

\begin{figure}[htbp]
  \centering
  \includegraphics[width=\columnwidth]{d1_growth_index}
  \caption{Normalised growth index (base = 100 at each language's first
    observed month). TypeScript reaches an index value of $\sim$20,000$\times$
    its 2015 baseline by 2025, dwarfing all other languages. Go and Rust
    (grey/unlabelled lines) also show multi-hundred-percent growth, while
    JavaScript, Python, Java, and PHP remain near their baselines, reflecting
    absolute-scale dominance rather than proportional growth.}
  \label{fig:d1_growth}
\end{figure}

\subsection{D2: 24-Month Holt-Winters Forecasts}

Forecasts were generated for all \textbf{30} top-ranked languages (minimum
36-month history required; all 30 met this criterion), producing 4,500 total
output rows (150 per language: 126 historical + 24 forecast).

Projected 24-month changes from June 2025 baseline:
\begin{itemize}
  \item \textbf{+27\%}: Nix — fastest projected growth, driven by
        reproducible build system adoption.
  \item \textbf{+26\%}: C\# — strong growth in game development (Unity)
        and cross-platform (.NET MAUI) segments.
  \item \textbf{+18\%}: Go, Rust, Julia — continued systems/ML momentum.
  \item \textbf{+10\%}: TypeScript, Python — sustained expansion at scale.
  \item \textbf{+4\%}: JavaScript — plateau indicating market saturation.
  \item \textbf{$-6$ to $-27$\%}: Vim Script, SCSS, Objective-C, Ruby, Lua,
        PHP — consistent multi-year decline captured by the trend state.
\end{itemize}

\begin{figure}[htbp]
  \centering
  \includegraphics[width=\columnwidth]{d2_top6_forecasts}
  \caption{24-month Holt-Winters ETS forecasts for the six highest-growth
    languages (blue = historical actual, orange dashed = forecast). All six
    show upward-trending projections, with the damped component preventing
    unrealistic straight-line extrapolation at the horizon.}
  \label{fig:d2}
\end{figure}

\begin{figure}[htbp]
  \centering
  \includegraphics[width=\columnwidth]{d2_forecast_summary}
  \caption{Projected 24-month percentage change for all 30 forecasted
    languages (green = growth, red = decline). Nix (+27\%), C\# (+26\%),
    Go and Rust (+18\%) lead; Vim Script ($-$27\%) and SCSS show the
    steepest projected declines.}
  \label{fig:d2_summary}
\end{figure}

\subsection{D3: Technology Migration Graph}

After allowlist filtering and minimum edge-weight thresholding, the migration
graph has \textbf{16 nodes} and \textbf{11 edges}. Top destination technologies
by total migration flow: CircleCI ($n{=}24$), Jest ($n{=}18$), Go ($n{=}12$).
This indicates active migration into modern CI/CD (CircleCI replacing Jenkins
or Travis) and testing (Jest replacing Mocha or Jasmine).

\begin{figure}[htbp]
  \centering
  \includegraphics[width=\columnwidth]{d3_network}
  \caption{Technology migration graph (16 nodes, 11 edges). Node size
    encodes PageRank; colour encodes net flow (green = net destination,
    red = net source). Jest and Go are the strongest migration attractors;
    Mocha and Travis are the primary migration sources.}
  \label{fig:d3}
\end{figure}

\subsection{D4: Developer Influence (PageRank + Betweenness)}

The developer co-contribution graph covers \textbf{25,000} developers with:
mean degree = 7.7, max degree = 211, isolated nodes (degree = 0) = 9,625
(38.5\%). The high isolated-node fraction reflects the reality that even
after bot filtering, many active developers contribute exclusively to one or
two projects and never co-occur with others in the top-25k selection.

PageRank ($\alpha = 0.85$, weighted by shared repositories) identifies the
most influential connectors in the open-source network — developers whose
removal would most degrade network reachability. Betweenness centrality
complements this by identifying bridge nodes that connect otherwise-disjoint
collaboration clusters.

\begin{figure}[htbp]
  \centering
  \includegraphics[width=\columnwidth]{d4_pagerank_vs_commits}
  \caption{PageRank vs.\ log total commits for all 25,000 developers
    ($r = 0.277$). The weak positive correlation confirms that high commit
    count alone does not determine network influence — many high-commit
    developers (bots) have near-zero PageRank while well-connected
    cross-project contributors rank highly.}
  \label{fig:d4}
\end{figure}

\subsection{D5: Community Detection}

Leiden community detection yields \textbf{10,255 communities} from 25,000
developers. Community labels were built by sampling 50 random Parquet shards
from \texttt{raw/developers/} (19.2M dev-repo edges), since loading the full
14\,GB dataset into Jupyter RAM caused an out-of-memory crash. The distribution
is heavily right-skewed:

\begin{itemize}
  \item Largest community: 1,679 members (a dense cluster of web ecosystem
        contributors sharing many JavaScript/TypeScript repositories).
  \item Median community size: 1 member (singleton).
  \item Singleton communities: 9,625 (93.9\% of all communities).
\end{itemize}

This heavy-tailed structure is characteristic of open-source collaboration
networks~\cite{borges2016popularity}: a small number of large, actively
maintained projects (React, Kubernetes, VSCode) attract dense contributor
clusters, while the majority of open-source work happens in small, isolated
projects or single-person repositories.

\begin{figure}[htbp]
  \centering
  \includegraphics[width=\columnwidth]{d5_community_sizes}
  \caption{Leiden community structure. \textit{Left}: size distribution
    (log scale) showing a heavy-tailed power-law pattern — 93.9\% of
    communities are singletons. \textit{Right}: top-20 communities by member
    count, labelled by their top contributing organisations. The largest
    community (1,679 members) is dominated by web-ecosystem contributors.}
  \label{fig:d5}
\end{figure}

% ─────────────────────────────────────────────────────────────────────────────
\newpage
\FloatBarrier
\section{Design Decisions Summary}

Table~\ref{tab:decisions} summarises the major design choices made during
implementation, including the alternatives considered and the reason each
was rejected.

\begin{table*}[!t]
\caption{Key Design Decisions and Trade-offs}
\label{tab:decisions}
\centering\small
\begin{tabular}{p{2.4cm} p{3.0cm} p{3.0cm} p{6.8cm}}
\toprule
\textbf{Decision Point} & \textbf{Chosen} & \textbf{Alternative} & \textbf{Reason} \\
\midrule
Language data source & HuggingFace 40M (2025) & BigQuery snapshot (2017) &
  BQ snapshot is point-in-time; HF provides creation timestamps and 8 more years of coverage \\
\addlinespace
Commit data source & GH Archive PushEvents & BigQuery frozen commits &
  BQ commits table is frozen at 2017; GH Archive is live and updated monthly \\
\addlinespace
Forecasting model & Holt-Winters ETS (damped) & ARIMA(2,1,2), Prophet &
  ARIMA projects flat lines for multi-year growth; Prophet requires CmdStan which fails on Dataproc \\
\addlinespace
Developer selection & Distinct repository count & Total commit count &
  Total commits selects CI bots; distinct repos selects real cross-project contributors \\
\addlinespace
Community detection & Leiden & Louvain &
  Leiden guarantees well-connected communities; Louvain can produce internally disconnected partitions \\
\addlinespace
Influence metric & PageRank + Betweenness & PageRank + HITS &
  HITS hub $\equiv$ authority on undirected graphs; betweenness identifies structural bridge nodes \\
\addlinespace
Storage format & Parquet + Snappy & CSV, ORC, AVRO &
  Columnar reads + predicate pushdown; Snappy gives fast decompression for compute-bound jobs \\
\addlinespace
Megarepo exclusion & \texttt{MAX\_DEVS\_}\newline\texttt{PER\_REPO=200} & No exclusion &
  Linux/TensorFlow contributor lists cause $O(N^2)$ join explosion: 50M+ edge rows from one repo \\
\addlinespace
Job orchestration & Sequential + wait-poll & Parallel (\texttt{-{}-async}) &
  Parallel jobs corrupt GCS outputs when a downstream job reads a path still being written \\
\addlinespace
Data cutoff & 2025-06 & 2025-12 or no cap &
  HF dataset ends July 23 2025; post-cutoff months show 49\% artificial repo-count drop \\
\bottomrule
\end{tabular}
\end{table*}

% ─────────────────────────────────────────────────────────────────────────────
\section{Conclusion \& Next Steps}

The full pipeline is operational end-to-end and publicly available~\cite{github_repo}. All six Spark jobs run
successfully on Dataproc (us-east1), and five analytical notebooks produce
quantitative outputs. The main findings to date: TypeScript and Rust exhibit
the strongest sustained growth trajectories; CircleCI and Jest are the
primary migration attractors; developer collaboration follows a heavy-tailed
community structure; and Nix, C\#, and Go are projected to lead adoption
growth over the next 24 months.

Remaining work for the final phase:
\begin{itemize}
  \item Enrich the D3 migration graph with PR/issue text and CHANGELOG
        parsing to increase edge coverage beyond the current 11 edges.
  \item Add Holt-Winters confidence intervals to the D2 notebook.
  \item Cross-join D4 developer PageRank with D1 language clusters to
        determine which developer communities are driving adoption of
        fast-growing languages.
  \item Evaluate whether D5 community language composition aligns with
        D1 language cluster boundaries.
  \item Develop an interactive web UI to expose the pipeline outputs —
        allowing users to explore language trend charts, browse forecasts,
        inspect the migration graph, and query developer community
        memberships through a dashboard interface.
\end{itemize}

% ─────────────────────────────────────────────────────────────────────────────
\section*{Acknowledgments}
The authors thank the Gebze Technical University Department of Computer
Engineering for support, and acknowledge Google Cloud Platform for cloud
computing credits used for this project.

% ─────────────────────────────────────────────────────────────────────────────
\clearpage  % flush all deferred wide floats (figure*, table*) before references
\begin{thebibliography}{00}

\bibitem{dyer2013boa}
R.\ Dyer, H.\ A.\ Nguyen, H.\ Rajan, and T.\ N.\ Nguyen,
``Boa: A language and infrastructure for analyzing ultra-large-scale software
repositories,''
in \textit{Proc.\ 35th Int.\ Conf.\ Software Engineering (ICSE)}, 2013,
pp.\ 422--431.

\bibitem{borges2016popularity}
H.\ Borges, A.\ Hora, and M.\ T.\ Valente,
``Understanding the factors that impact the popularity of GitHub repositories,''
in \textit{Proc.\ IEEE Int.\ Conf.\ Software Maintenance and Evolution (ICSME)},
2016, pp.\ 334--344.

\bibitem{bissyande2013popularity}
T.\ F.\ Bissyandé, D.\ Lo, L.\ Jiang, L.\ Réveillère, J.\ Klein, and Y.\ Le Traon,
``Got issues? who cares about it? a large scale investigation of issue trackers
from GitHub,''
in \textit{Proc.\ IEEE Int.\ Symp.\ Software Reliability Engineering (ISSRE)},
2013, pp.\ 188--197.

\bibitem{gardner1985exponential}
E.\ S.\ Gardner,
``Exponential smoothing: The state of the art,''
\textit{Journal of Forecasting}, vol.\ 4, no.\ 1, pp.\ 1--28, 1985.

\bibitem{makridakis2018m4}
S.\ Makridakis, E.\ Spiliotis, and V.\ Assimakopoulos,
``The M4 competition: Results, findings, conclusion and way forward,''
\textit{International Journal of Forecasting}, vol.\ 34, no.\ 4,
pp.\ 802--808, 2018.

\bibitem{hyndman2018forecasting}
R.\ J.\ Hyndman and G.\ Athanasopoulos,
\textit{Forecasting: Principles and Practice}, 2nd ed.
Melbourne: OTexts, 2018. [Online]. Available: \url{https://otexts.com/fpp2/}

\bibitem{traag2019leiden}
V.\ A.\ Traag, L.\ Waltman, and N.\ J.\ van Eck,
``From Louvain to Leiden: Guaranteeing well-connected communities,''
\textit{Scientific Reports}, vol.\ 9, no.\ 1, p.\ 5233, 2019.

\bibitem{brin1998pagerank}
S.\ Brin and L.\ Page,
``The anatomy of a large-scale hypertextual web search engine,''
\textit{Computer Networks and ISDN Systems}, vol.\ 30, no.\ 1--7,
pp.\ 107--117, 1998.

\bibitem{github_repo}
S.\ Bardakc\i, B.\ D.\ G\"{o}ksu, and Z.\ Yi\u{g}it,
``BigData-Github\_Analysis,'' GitHub, 2025.
[Online]. Available: \url{https://github.com/selin-bardakci/BigData-Github_Analysis}

\end{thebibliography}

\end{document}
